%=============================================================================!
% 6.S  WHAT WE NOW HAVE                                                       !
%=============================================================================!
\unnumbered{What we now have}

Mechanics can be written in any coordinates that fix the positions of a
system once its constraints are obeyed, its generalised coordinates, with
the Lagrangian $\Lag = K - U$ expressed in them. Among all paths with the
same ends in space and time, the true one makes the action $\Act = \int
\Lag\,\dd t$ stationary, which is not always least. Changing a path by a
small bump and integrating by parts turns this into the Euler--Lagrange
equations, $\dd(\partial\Lag/\partial\dot{q}_k)/\dd t = \partial\Lag/
\partial q_k$, one for each coordinate, which for forces from a potential
are Newton's second law.

Constraints that do no work drop out of $\Lag$:
solving the bead's equation along its wire gives the times of the journey
that Chapter~3 described by its speeds, and the pendulum's equation takes
one line.

The momentum $p_k = \partial\Lag/\partial\dot{q}_k$ of a
coordinate that $\Lag$ does not contain is conserved, which is why a
central force conserves angular momentum; the distance from the centre
then moves in the effective potential $U + L^2/(2mr^2)$.

When a
constraining force is wanted, or no convenient coordinates exist, a Lagrange multiplier keeps the constraint: for a pendulum it gives the
force of the rod, $-\lambda L$, and for bonds held at fixed length the
multipliers are what Chapter~11 computes at every step.

Chapter~7 trades the generalised velocities for the momenta and arrives at
Hamilton's equations, the form of mechanics on which the methods of
simulation are built.
